\uuid{jnUW}
\titre{Intégration par parties/linéarité}
\auteur{hanani ; mzari}
\organisation{exo7}
\licence{CC BY-NC-SA 4.0}
\datecreate{2020-01-20}
\niveau{L1}
\module{Analyse}
\chapitre{Calcul d'intégrales}
\sousChapitre{Autre}
\difficulte{3}
\theme{Calculs d'intégrales, intégration par parties, linéarité}
\exo7qcmid{lille-2.integrales.021}

\contenu{
\begin{question}
\qclassification{127.99}
\qtags{motcle=intégration par parties/linéarité}
On pose $\displaystyle I=\int _0^{\pi/2}x\cos ^2x\, \mathrm{d}x$, $\displaystyle J=\int _0^{\pi/2}x\sin ^2x\, \mathrm{d}x$ et $\displaystyle K=\int _0^{\pi/2}x\cos (2x)\, \mathrm{d}x$. Parmi les affirmations suivantes, cocher celles qui sont vraies :
\begin{answers}
\bad{Une intégration par parties donne : $\displaystyle K=\frac{1}{2}$.}
\good{$\displaystyle I+J=\frac{\pi ^2}{8}$ et $I-J=K$.}
\bad{$\displaystyle I=\frac{\pi ^2+4}{16}$.}
\bad{$\displaystyle J=\frac{\pi ^2-4}{16}$.}
\end{answers}
\vskip2mm
\begin{explanations}
Une intégration par parties, avec $u=x$ et $v=1/2\sin (2x)$, donne
$$K=\left[\frac{x}{2}\sin (2x)\right]_0^{\pi/2}-\frac{1}{2}\int _0^{\pi/2} \sin (2x)\, \mathrm{d}x=-\frac{1}{2}.$$
A l'aide des relations $\cos ^2x+\sin ^2x=1$ et  $\cos ^2x-\sin ^2x=\cos (2x)$, on obtient :
$$I+J=\int _0^{\pi/2}x\mathrm{d}x=\frac{\pi ^2}{8}\mbox{ et }I-J=K=-\frac{1}{2}.$$
La somme et la différence de ces égalités donnent : $\displaystyle I=\frac{\pi ^2-4}{16}$ et $\displaystyle J=\frac{\pi ^2+4}{16}$.
\end{explanations}
\end{question}
}
